% =============================================================
%  Chapitre 3 — Réalisation
% =============================================================
\chapter{Réalisation}

\section{Introduction}

Ce chapitre présente la réalisation concrète de l'application Telnet Test
Manager. Nous parcourons successivement la structure du projet backend,
l'implémentation du moteur Telnet, l'API REST, le modèle de données
MongoDB, les interfaces utilisateur, la conteneurisation Docker, le
déploiement Kubernetes/Helm, le pipeline GitHub Actions et la configuration
du runner auto-hébergé.

% -----------------------------------------------------------------
\section{Implémentation}

\subsection{Structure du projet Node.js/Express}

Le backend adopte la \textbf{Clean Architecture} en quatre couches dont
les dépendances ne pointent que vers l'intérieur (Domain ← Application
← Infrastructure ← Interfaces).

\begin{lstlisting}[language={},caption={Arborescence du projet backend},label={lst:arbo_backend}]
backend/
├── src/
│   ├── domain/                   # Entités et interfaces (aucune dépendance externe)
│   │   ├── entities/             # User, Slot, TelnetCommand, TestResult...
│   │   └── repositories/        # IUserRepository, ISlotRepository...
│   ├── application/
│   │   └── usecases/             # LoginUseCase, RunTestUseCase, GetResultsUseCase...
│   ├── infrastructure/
│   │   ├── config/database.js    # Connexion MongoDB
│   │   ├── database/
│   │   │   ├── models/           # 9 schémas Mongoose
│   │   │   └── repositories/    # 9 implémentations MongoDB
│   │   ├── telnet/
│   │   │   ├── TestWorkerManager.js  # Gestion du cycle de vie des Workers
│   │   │   └── testWorker.js     # Logique Telnet isolée dans un Worker Thread
│   │   └── metrics/
│   │       └── prometheusMetrics.js  # 4 métriques Prometheus
│   ├── interfaces/
│   │   ├── http/
│   │   │   ├── controllers/      # 9 contrôleurs minces
│   │   │   ├── routes/           # 10 groupes de routes Express
│   │   │   └── middlewares/      # authenticate, requireRole, auditLog, errorHandler
│   │   └── websocket/
│   │       └── WebSocketServer.js  # Serveur WS port 3003
│   └── main/
│       ├── server.js             # Point d'entrée
│       ├── app.js                # Configuration Express
│       └── container.js         # Conteneur d'injection de dépendances
├── Dockerfile
└── package.json
\end{lstlisting}

\subsection{Module de test Telnet}

Le cœur du système est le moteur Telnet, composé du
\textbf{TestWorkerManager} (gestionnaire du cycle de vie) et du
\textbf{testWorker} (logique d'exécution dans un Worker Thread isolé).

\begin{lstlisting}[language=JavaScript,
  caption={testWorker.js — exécution d'un test Telnet en mode séquence},
  label={lst:telnet_worker}]
const { workerData, parentPort } = require('worker_threads');
const Telnet = require('telnet-client');

async function runTest() {
  const { slot, command, testId } = workerData;
  const telnet = new Telnet();

  // Paramètres de connexion à l'équipement réseau
  const params = {
    host:                 slot.adresse,  // ex : 192.168.5.1
    port:                 slot.port,     // 23 (Telnet)
    timeout:              15000,
    shellPrompt:          /#\s*$/,       // Invite shell Linux (root)
    loginPrompt:          /login:\s*$/,
    passwordPrompt:       /Password:\s*$/,
    username:             'root',
    password:             'root',
    negotiationMandatory: true,
    stripControls:        true,
    echoLines:            0,
  };

  try {
    await telnet.connect(params);

    if (command.type === 'single') {
      // ── Mode commande unitaire ─────────────────────────────
      const response = await telnet.exec(command.command);
      const success  = !command.expectedResponse ||
                       response.includes(command.expectedResponse);
      parentPort.postMessage({ type: 'completed', testId, success, response });

    } else if (command.type === 'sequence') {
      // ── Mode séquence : étapes ordonnées ──────────────────
      for (let i = 0; i < command.steps.length; i++) {
        const step = command.steps[i];
        const res  = await telnet.exec(step.command);
        const ok   = !step.expectedResponse ||
                     res.includes(step.expectedResponse);
        // Envoyer le résultat de l'étape au thread principal
        parentPort.postMessage({
          type: 'step', testId,
          stepIndex: i,
          status:    ok ? 'SUCCESS' : 'FAIL',
          response:  res,
          timestamp: new Date().toISOString(),
        });
      }
      parentPort.postMessage({ type: 'completed', testId, success: true });

    } else if (command.type === 'monitoring') {
      // ── Mode supervision : écoute d'événements ────────────
      await telnet.exec(command.command);
      const duration = command.steps[0]?.duration || 15000;
      await new Promise(resolve => setTimeout(resolve, duration));
      parentPort.postMessage({ type: 'completed', testId, success: true });
    }

    await telnet.end();

  } catch (err) {
    // Signaler l'erreur au thread principal
    parentPort.postMessage({ type: 'error', testId, message: err.message });
  }
}

runTest();
\end{lstlisting}

\begin{figure}[H]
  \centering
  \fbox{\parbox{0.82\textwidth}{\centering\vspace{0.8cm}
    {\large\color{gray!55}\textbf{Figure 3.1 — Résultat d'un test Telnet en temps réel}}\\[0.4cm]
    {\small Capture de l'interface Dashboard~: terminal de résultats avec statuts
     SUCCESS/FAIL par étape}\\[0.3cm]
    {\scriptsize\itshape\color{gray}(capture d'écran à insérer)}
    \vspace{0.8cm}}}
  \caption{Interface Dashboard~: résultat en temps réel d'un test Telnet}
  \label{fig:resultat_telnet}
\end{figure}

\subsection{API REST}

\begin{longtable}{|C{1.8cm}|L{5cm}|L{3cm}|C{2.5cm}|}
  \caption{Principaux endpoints de l'API REST}
  \label{tab:api_rest} \\
  \hline
  \textbf{Méthode} & \textbf{Endpoint} & \textbf{Description} & \textbf{Auth} \\
  \hline
  \endfirsthead
  \multicolumn{4}{|c|}{\tablename\ \thetable{} — suite} \\
  \hline
  \textbf{Méthode} & \textbf{Endpoint} & \textbf{Description} & \textbf{Auth} \\
  \hline
  \endhead
  \hline
  \endfoot
  \hline
  \endlastfoot

  POST   & \texttt{/login}              & Authentification JWT     & Non \\ \hline
  POST   & \texttt{/logout}             & Déconnexion              & JWT \\ \hline
  GET    & \texttt{/postes}             & Liste des postes         & JWT \\ \hline
  POST   & \texttt{/postes}             & Créer un poste           & Ingénieur \\ \hline
  GET    & \texttt{/produits}           & Liste des produits       & JWT \\ \hline
  GET    & \texttt{/slots}              & Liste des slots          & JWT \\ \hline
  POST   & \texttt{/slots}              & Créer un slot            & Ingénieur \\ \hline
  GET    & \texttt{/telnet-commands}    & Bibliothèque commandes   & JWT \\ \hline
  POST   & \texttt{/telnet-commands}    & Créer une commande       & Ingénieur \\ \hline
  POST   & \texttt{/run-test}           & Lancer un test Telnet    & JWT \\ \hline
  POST   & \texttt{/stop-test}          & Arrêter un test          & JWT \\ \hline
  POST   & \texttt{/stop-monitoring}    & Arrêter la supervision   & JWT \\ \hline
  GET    & \texttt{/test-results}       & Historique des résultats & JWT \\ \hline
  GET    & \texttt{/reports}            & Liste des rapports       & JWT \\ \hline
  POST   & \texttt{/reports/generate}   & Générer un rapport       & JWT \\ \hline
  GET    & \texttt{/admin/users}        & Gestion utilisateurs     & Admin \\ \hline
  POST   & \texttt{/admin/users}        & Créer un utilisateur     & Admin \\ \hline
  GET    & \texttt{/admin/audit-logs}   & Journal d'audit          & Admin \\ \hline
  GET    & \texttt{/admin/analytics}    & Statistiques globales    & Admin \\ \hline
  GET    & \texttt{/metrics}            & Métriques Prometheus     & Public \\ \hline
  GET    & \texttt{/health}             & Statut du serveur        & Public \\ \hline
\end{longtable}

\subsection{Modèle de données MongoDB}

\begin{lstlisting}[language=JavaScript,
  caption={Schéma Mongoose — collection \texttt{slots}},
  label={lst:schema_slot}]
const mongoose = require('mongoose');

const SlotSchema = new mongoose.Schema({
  id:          { type: Number, required: true, unique: true },
  nom:         { type: String, required: true },
  produitId:   { type: Number, required: true, ref: 'Produit' },
  adresse:     { type: String, required: true }, // Adresse IP : ex. "192.168.5.1"
  port:        { type: Number, default: 23 },    // Port Telnet standard
  description: { type: String, default: '' },
});

module.exports = mongoose.model('Slot', SlotSchema);
\end{lstlisting}

\begin{lstlisting}[language=JavaScript,
  caption={Schéma Mongoose — collection \texttt{testResults}},
  label={lst:schema_result}]
const TestResultSchema = new mongoose.Schema({
  id:        { type: Number, unique: true },   // Timestamp : Date.now()
  slotId:    { type: Number, ref: 'Slot' },
  posteId:   { type: Number, ref: 'Poste' },
  produitId: { type: Number, ref: 'Produit' },
  commandId: { type: String, ref: 'TelnetCommand' },
  runMode:   {
    type: String,
    enum: ['single', 'builtin_sequence', 'sequence', 'monitoring']
  },
  statut: {
    type: String,
    enum: ['PENDING', 'SUCCESS', 'FAIL', 'STOPPED'],
    default: 'PENDING'
  },
  startTime: { type: Date },
  endTime:   { type: Date },
  steps: [{
    step:        String,
    description: String,
    statut:      { type: String, enum: ['SUCCESS', 'FAIL'] },
    timestamp:   Date,
  }],
  logs: [{ type: String }],
});
\end{lstlisting}

\subsection{Interface utilisateur}

\begin{figure}[H]
  \centering
  \fbox{\parbox{0.88\textwidth}{\centering\vspace{1cm}
    {\large\color{gray!55}\textbf{Figure 3.2 — Tableau de bord principal (Dashboard)}}\\[0.4cm]
    {\small Panneau de sélection (Poste / Produit / Slot / Commande) + terminal
     de résultats en temps réel via WebSocket}\\[0.3cm]
    {\scriptsize\itshape\color{gray}(capture d'écran à insérer)}
    \vspace{1cm}}}
  \caption{Tableau de bord principal de l'application}
  \label{fig:dashboard}
\end{figure}

\begin{figure}[H]
  \centering
  \fbox{\parbox{0.88\textwidth}{\centering\vspace{1cm}
    {\large\color{gray!55}\textbf{Figure 3.3 — Formulaire de lancement d'un test Telnet}}\\[0.4cm]
    {\small Sélection du slot, choix de la commande, options d'exécution
     et bouton «~Lancer le test~»}\\[0.3cm]
    {\scriptsize\itshape\color{gray}(capture d'écran à insérer)}
    \vspace{1cm}}}
  \caption{Formulaire de lancement d'un test Telnet}
  \label{fig:form_test}
\end{figure}

\subsection{Conteneurisation Docker}

\begin{lstlisting}[language=docker,
  caption={backend/Dockerfile — image Node.js 20 Alpine},
  label={lst:dockerfile_back}]
# Image de base légère : Node.js 20 sur Alpine Linux (~50 MB)
FROM node:20-alpine

WORKDIR /app

# Copie des manifestes pour exploiter le cache Docker
COPY package*.json ./

# Installation des dépendances de production uniquement
RUN npm ci --only=production

# Cache-busting : force la recopie du code source à chaque build
ARG CACHEBUST=1
COPY . .

EXPOSE 3002

CMD ["node", "src/main/server.js"]
\end{lstlisting}

\begin{lstlisting}[language=docker,
  caption={frontend/Dockerfile — build multi-stage React + Nginx},
  label={lst:dockerfile_front}]
# ── Étape 1 : Construction de l'application React ─────────────────
FROM node:18-alpine AS builder
WORKDIR /app
COPY package*.json ./
RUN npm ci
COPY . .
# Génère les fichiers statiques optimisés dans ./build
RUN npm run build

# ── Étape 2 : Serveur Nginx pour les fichiers statiques ───────────
FROM nginx:alpine
# L'image finale ne contient que Nginx + les fichiers compilés (~25 MB)
COPY --from=builder /app/build /usr/share/nginx/html
COPY nginx.conf /etc/nginx/conf.d/default.conf
EXPOSE 80
CMD ["nginx", "-g", "daemon off;"]
\end{lstlisting}

\begin{lstlisting}[language=yaml,
  caption={docker-compose.yml — stack locale de développement (5 services)},
  label={lst:docker_compose}]
services:
  mongodb:
    image: mongo:7
    ports:
      - "27017:27017"
    environment:
      MONGO_INITDB_ROOT_USERNAME: admin
      MONGO_INITDB_ROOT_PASSWORD: <MOT_DE_PASSE>   # à définir dans .env
      MONGO_INITDB_DATABASE:      telnet-db
    volumes:
      - mongo-data:/data/db

  backend:
    build: ./backend
    ports:
      - "3002:3002"    # API REST
      - "3003:3003"    # WebSocket
    environment:
      MONGO_URI:      mongodb://admin:<MOT_DE_PASSE>@mongodb:27017/telnet-db?authSource=admin
      JWT_SECRET:     ${JWT_SECRET}   # injecté depuis le fichier .env
      ALLOWED_ORIGIN: http://localhost:3000
    depends_on: [mongodb]

  frontend:
    build: ./frontend
    ports:
      - "3000:80"
    depends_on: [backend]

  prometheus:
    image: prom/prometheus:latest
    ports:
      - "9090:9090"
    volumes:
      - ./monitoring/prometheus.yml:/etc/prometheus/prometheus.yml

  grafana:
    image: grafana/grafana:latest
    ports:
      - "3001:3000"
    environment:
      GF_SECURITY_ADMIN_PASSWORD: <MOT_DE_PASSE>   # à définir dans .env
    volumes:
      - ./monitoring/grafana-datasources.yaml:/etc/grafana/provisioning/datasources/ds.yaml
      - ./monitoring/grafana-dashboard.json:/etc/grafana/provisioning/dashboards/telnet.json
    depends_on: [prometheus]

volumes:
  mongo-data:
\end{lstlisting}

\subsection{Déploiement Kubernetes avec Helm}

\begin{lstlisting}[language={},
  caption={Structure du chart Helm \texttt{telnet-app} (v0.3.0)},
  label={lst:helm_structure}]
helm/telnet-app/
├── Chart.yaml                    # Métadonnées (name, version: 0.3.0)
├── values.yaml                   # 40+ valeurs paramétrables
└── templates/
    ├── backend-deployment.yaml   # Deployment backend (2 répliques)
    ├── backend-hpa.yaml          # HPA : 1-5 répliques, seuil CPU 70 %
    ├── backend-pvc.yaml          # PVC 1 Gi pour /app/reports
    ├── backend-service.yaml      # ClusterIP : 3002
    ├── frontend-deployment.yaml  # Deployment frontend (1 réplique)
    ├── frontend-service.yaml     # ClusterIP : 80
    ├── mongodb-deployment.yaml   # Deployment MongoDB 7
    ├── mongodb-pvc.yaml          # PVC 2 Gi pour /data/db
    ├── mongodb-service.yaml      # ClusterIP : 27017
    ├── configmap.yaml            # Variables d'environnement
    ├── secret.yaml               # JWT_SECRET, credentials MongoDB
    ├── ingress.yaml              # Nginx ingress : / → frontend, /api → backend
    └── seed-job.yaml             # Hook post-install : seed admin + engineer
\end{lstlisting}

\begin{lstlisting}[language=yaml,
  caption={Extrait — backend-deployment.yaml (chart Helm)},
  label={lst:k8s_deploy}]
apiVersion: apps/v1
kind: Deployment
metadata:
  name: {{ .Release.Name }}-backend
  namespace: {{ .Release.Namespace }}
spec:
  replicas: {{ .Values.replicaCount.backend }}
  strategy:
    type: RollingUpdate
    rollingUpdate:
      maxUnavailable: 0   # Aucun pod indisponible pendant la mise à jour
      maxSurge:       1   # Un pod supplémentaire temporaire autorisé
  selector:
    matchLabels:
      app: {{ .Release.Name }}-backend
  template:
    spec:
      # Attendre que MongoDB soit prêt avant de démarrer le backend
      initContainers:
        - name: wait-for-mongodb
          image: busybox:1.35
          command: ['sh', '-c',
            'until nc -z {{ .Release.Name }}-mongodb 27017;
             do echo waiting; sleep 2; done']
      containers:
        - name: backend
          image: "telnet-backend:{{ .Values.image.backend.tag }}"
          ports:
            - containerPort: 3002  # API REST
            - containerPort: 3003  # WebSocket
          resources:
            requests:
              memory: {{ .Values.resources.backend.requests.memory }}
              cpu:    {{ .Values.resources.backend.requests.cpu }}
            limits:
              memory: {{ .Values.resources.backend.limits.memory }}
              cpu:    {{ .Values.resources.backend.limits.cpu }}
\end{lstlisting}

\begin{lstlisting}[language=yaml,
  caption={backend-hpa.yaml — Horizontal Pod Autoscaler},
  label={lst:hpa}]
apiVersion: autoscaling/v2
kind: HorizontalPodAutoscaler
metadata:
  name: {{ .Release.Name }}-backend-hpa
spec:
  scaleTargetRef:
    apiVersion: apps/v1
    kind:       Deployment
    name:       {{ .Release.Name }}-backend
  minReplicas: {{ .Values.hpa.backend.minReplicas }}   # 1
  maxReplicas: {{ .Values.hpa.backend.maxReplicas }}   # 5
  metrics:
    - type: Resource
      resource:
        name: cpu
        target:
          type:               AverageUtilization
          averageUtilization: {{ .Values.hpa.backend.targetCPUUtilizationPercentage }}  # 70
\end{lstlisting}

\subsection{Pipeline GitHub Actions}

\begin{lstlisting}[language=yaml,
  caption={.github/workflows/ci-cd.yaml — pipeline CI/CD complet},
  label={lst:pipeline}]
name: "CI/CD — Telnet Test Manager"

on:
  push:
    branches: [main, develop]
  pull_request:
    branches: [main]

# Annuler les exécutions précédentes pour la même branche
concurrency:
  group: deploy-${{ github.ref }}
  cancel-in-progress: true

jobs:
  # ────────────────────────────────────────────────────────────
  # JOB 1 — Intégration continue Backend
  # ────────────────────────────────────────────────────────────
  ci-backend:
    runs-on: ubuntu-latest
    steps:
      - uses: actions/checkout@v4
      - uses: actions/setup-node@v4
        with: { node-version: '20' }
      - run: cd backend && npm ci
      - run: cd backend && npm audit --audit-level=critical
      - run: cd backend && npm test --if-present

  # ────────────────────────────────────────────────────────────
  # JOB 2 — Intégration continue Frontend
  # ────────────────────────────────────────────────────────────
  ci-frontend:
    runs-on: ubuntu-latest
    steps:
      - uses: actions/checkout@v4
      - uses: actions/setup-node@v4
        with: { node-version: '20' }
      - run: cd frontend && npm ci
      - run: cd frontend && npm audit --audit-level=critical
      - run: cd frontend && npm run build
        env: { CI: false }

  # ────────────────────────────────────────────────────────────
  # JOB 3 — Déploiement (self-hosted, branche main uniquement)
  # ────────────────────────────────────────────────────────────
  deploy:
    runs-on: self-hosted
    needs: [ci-backend, ci-frontend]
    if: github.ref == 'refs/heads/main'
    steps:
      - uses: actions/checkout@v4

      - name: "Démarrer Minikube si nécessaire"
        run: minikube status || minikube start --driver=docker

      - name: "Construire les images Docker (daemon Minikube)"
        run: |
          eval $(minikube docker-env)
          docker build -t telnet-backend:latest  ./backend
          docker build -t telnet-frontend:latest ./frontend

      - name: "Déploiement Helm (atomic)"
        run: |
          helm upgrade --install telnet-app ./helm/telnet-app \
            --create-namespace --namespace telnet-app \
            --set secrets.mongoRootPassword=${{ secrets.MONGO_PASSWORD }} \
            --set secrets.jwtSecret=${{ secrets.JWT_SECRET }}            \
            --atomic --timeout 600s

      - name: "Attendre la disponibilité des pods"
        run: |
          kubectl rollout status deployment/telnet-app-backend  \
            -n telnet-app --timeout=240s
          kubectl rollout status deployment/telnet-app-frontend \
            -n telnet-app --timeout=180s

      - name: "Résumé du déploiement"
        run: kubectl get pods,services -n telnet-app
\end{lstlisting}

\begin{figure}[H]
  \centering
  \fbox{\parbox{0.88\textwidth}{\centering\vspace{0.8cm}
    {\large\color{gray!55}\textbf{Figure 3.4 — Exécution du pipeline dans GitHub Actions}}\\[0.4cm]
    {\small Vue de l'onglet Actions~: 3 jobs (CI Backend, CI Frontend, Deploy)
     avec statuts Success/Failure}\\[0.3cm]
    {\scriptsize\itshape\color{gray}(capture d'écran à insérer)}
    \vspace{0.8cm}}}
  \caption{Exécution réussie du pipeline CI/CD dans GitHub Actions}
  \label{fig:pipeline_github}
\end{figure}

\subsection{Self-hosted runner}

Le runner auto-hébergé est installé sur la machine Windows locale.
L'enregistrement auprès du dépôt GitHub s'effectue via~:

\begin{lstlisting}[language={},
  caption={Commandes d'enregistrement du runner auto-hébergé},
  label={lst:runner_setup}]
# Depuis le dossier d'installation du runner GitHub Actions
.\config.cmd --url   https://github.com/[ORGANISATION]/[REPO] `
             --token [TOKEN_RUNNER]                           `
             --name  "minikube-runner-local"                  `
             --labels self-hosted,windows,minikube

# Démarrage du runner en écoute permanente
.\run.cmd
\end{lstlisting}

\begin{figure}[H]
  \centering
  \fbox{\parbox{0.82\textwidth}{\centering\vspace{0.8cm}
    {\large\color{gray!55}\textbf{Figure 3.5 — Runner auto-hébergé actif sur GitHub}}\\[0.4cm]
    {\small Paramètres du dépôt $\to$ Actions $\to$ Runners~: statut
     «~Idle~» (en attente de jobs)}\\[0.3cm]
    {\scriptsize\itshape\color{gray}(capture d'écran à insérer)}
    \vspace{0.8cm}}}
  \caption{Runner auto-hébergé enregistré et actif}
  \label{fig:runner_github}
\end{figure}

\section{Conclusion}

Ce chapitre a détaillé la réalisation complète de Telnet Test Manager,
depuis l'implémentation du moteur Telnet en Worker Threads jusqu'au
pipeline de déploiement GitHub Actions. La conteneurisation Docker
multi-stage, l'orchestration Kubernetes/Helm avec autoscaling horizontal
et la surveillance Prometheus/Grafana constituent une infrastructure
DevOps de niveau production. L'ensemble des fonctionnalités réalisées
répondent aux seize besoins fonctionnels identifiés au chapitre précédent.
